%! Author = markt
%! Date = 3/10/2025

% Preamble
\documentclass[11pt]{article}

% Packages
\usepackage{latexsym,amsfonts,amssymb,amsthm,amsmath, yfonts, cancel, listings, xcolor, graphicx}
\usepackage{color}
\usepackage{listings}

\author{Mark Stanley}

\definecolor{codegreen}{rgb}{0,0.6,0}
\definecolor{codegray}{rgb}{0.5,0.5,0.5}
\definecolor{codepurple}{rgb}{0.58,0,0.82}
\definecolor{backcolour}{rgb}{0.95,0.95,0.92}
\lstdefinestyle{mystyle}{
    backgroundcolor=\color{backcolour},
    commentstyle=\color{codegreen},
    keywordstyle=\color{magenta},
    numberstyle=\tiny\color{codegray},
    stringstyle=\color{codepurple},
    basicstyle=\ttfamily\footnotesize,
    breakatwhitespace=false,
    breaklines=true,
    captionpos=b,
    keepspaces=true,
    numbers=left,
    numbersep=5pt,
    showspaces=false,
    showstringspaces=false,
    showtabs=false,
    tabsize=2
}

\title{Math 443 Page Rank}

% Document
\begin{document}
    \lstset{style=mystyle}

    \maketitle

    \subsection{Page Rank}
    Consider all websites in a set as being a graph between each other $g_i \in G$
    where $G$ is the set of all $n$ websites.
    Each connection between two nodes is a link.
    Represent each node as a vector $v_i$, which is a one-hot vector of size $n$, where presence indicates a link to the website.
    We then normalize each vector $v_i$ such that all values add to $1$ (just softmax).
    Create a matrix $L$ of size $n \times n$ where each column is a vector $v_i$.
    Define $r_A = \sum^n_{j=1}L_{A,j}r_j$, which is the rank of the website $A$.
    This is just $r=Lr$, where $r$ is the rank vector.
    We perform this operation until convergence, which is when the rank vector does not change.

\end{document}